% =============================================================================
%  Weeks 5-7 Progress Report
%  Baseline probabilistic progression model
%
%  Figures come from outputs/model/figures, regenerated by
%      python scripts/03_progression_model.py
% =============================================================================
\documentclass[11pt,a4paper]{article}

\usepackage[T1]{fontenc}
\usepackage[utf8]{inputenc}
\usepackage[margin=2.4cm]{geometry}
\usepackage{booktabs}
\usepackage{longtable}
\usepackage{array}
\usepackage{xcolor}
\usepackage{graphicx}
\usepackage{enumitem}
\usepackage{caption}
\usepackage{amsmath}
\usepackage{amssymb}
\usepackage{float}
\usepackage[hidelinks]{hyperref}

\graphicspath{{../outputs/model/figures/}}

\definecolor{warn}{HTML}{B7770B}

\newcommand{\code}[1]{\texttt{\small #1}}
\newcommand{\arr}{$\rightarrow$}

\setlist[itemize]{itemsep=2pt,topsep=4pt}
\setlist[enumerate]{itemsep=2pt,topsep=4pt}
\captionsetup{font=small,labelfont=bf}

\newenvironment{keypoint}
  {\begin{center}\begin{minipage}{0.94\textwidth}%
   \vspace{4pt}\hrule\vspace{7pt}\small}
  {\vspace{7pt}\hrule\vspace{4pt}\end{minipage}\end{center}}

\title{\textbf{A Probabilistic Model of Motor Subtype Progression}\\[0.6em]
       \large Weeks 5--7 Progress Report}
\author{Parv\\[0.3em]
        \small Supervisor: Prof.\ K.~R.\ Sahasranand\\
        \small Department of Data Science, IIT Palakkad}
\date{August 2026}

\begin{document}
\maketitle

\begin{abstract}
\noindent
Weeks 5--7 called for a baseline probabilistic model of subtype transition,
validated on held-out patients. Three were built and compared: the discrete
Markov chain the proposal originally sketched, a continuous-time chain that
respects the actual gap between visits, and a hidden Markov model in which a
latent phenotype is observed through a misclassification matrix.

The two ordinary chains are both misspecified, and a direct test says why. In
the observed labels, the next state depends strongly on the \emph{previous}
one even after conditioning on the current one --- for patients currently
labelled TD, arriving from PIGD raises the chance of being PIGD next visit from
0.086 to 0.393 ($\chi^2 = 156.6$, $p = 7.9 \times 10^{-33}$). A Markov chain in
the label cannot represent that. A latent state read through noise can, and
does: the hidden Markov model beats the discrete baseline on held-out patients
in \textbf{10 of 10} independent splits, by $+0.028 \pm 0.009$ log-likelihood
per transition (paired $t$-test $p = 5.4 \times 10^{-6}$).

The fitted model delivers the output table the proposal specifies --- the
probability of occupying each subtype at 1, 2, 3 and 5 years --- and makes two
substantive claims: the underlying phenotype is far more stable than its label
(TD persists about 10 years latent against 3.4 observed), and progression is
\emph{sequential}, passing through Indeterminate rather than jumping from TD
to PIGD.
\end{abstract}

\vspace{0.3em}

% =============================================================================
\section{Objective}
% =============================================================================

The proposal assigns Weeks 5--7:

\begin{itemize}
  \item Develop a simple probabilistic model (Markov or transition-probability)
        from the observed longitudinal transitions.
  \item Estimate, for each current subtype, the probability of remaining stable
        or transitioning at the next follow-up.
  \item Validate the model using held-out patients.
\end{itemize}

\noindent
with the deliverable stated as an \emph{initial progression model providing
probabilities of subtype transition}. All three are complete. The model that
survives validation is not the one originally sketched, and Section~3 explains
why.

% =============================================================================
\section{Design}
% =============================================================================

\subsection{Data}

Input is \code{outputs/visit\_panel.csv} from Weeks 1--2: 4{,}441 labelled
visits over 1{,}080 sporadic PD patients, yielding 3{,}361 consecutive
transition pairs from the 934 patients with at least two labelled visits.

\subsection{Held-out split}

Splitting is at the level of the \emph{patient}, never the visit. A patient
contributes several transitions; splitting those independently would put the
same person on both sides and make the held-out estimate optimistic.

\begin{table}[htbp]
\centering
\begin{tabular}{@{}lrrrrr@{}}
\toprule
\textbf{Split} & \textbf{TD} & \textbf{Indet.} & \textbf{PIGD} & \textbf{Patients} & \textbf{Transitions} \\
\midrule
Train & 543 & 76 & 136 & 755 & 2{,}359 \\
Test  & 233 & 33 &  59 & 325 & 1{,}002 \\
\bottomrule
\end{tabular}
\caption{Patient-level split (70/30, seed 42), stratified by baseline subtype:
71.9\% TD on the training side against 71.7\% on the test side. The manifest is
written to \code{outputs/model/split\_manifest.csv}.}
\end{table}

\noindent
Of these, 658 training and 276 test patients have the two or more labelled
visits a sequence model requires.

% =============================================================================
\section{Why a Markov chain in the label is not enough}
% =============================================================================

A Markov chain assumes that the current state is a sufficient summary of the
past: $P(\text{next} \mid \text{current}, \text{previous}) =
P(\text{next} \mid \text{current})$. That is testable directly, on the 2{,}427
labelled triples in the panel.

\begin{table}[htbp]
\centering
\begin{tabular}{@{}lrrr@{}}
\toprule
\textbf{Currently} & \textbf{Triples} & \textbf{$\chi^2$ (dof 4)} & \textbf{$p$} \\
\midrule
TD            & 1{,}583 & 156.6 & $7.9\times10^{-33}$ \\
Indeterminate &    235 &  15.3 & $4.2\times10^{-3}$ \\
PIGD          &    609 &  61.2 & $1.7\times10^{-12}$ \\
\bottomrule
\end{tabular}
\caption{Test of the Markov property. Rejected in every state.}
\end{table}

\noindent
Concretely, for the 1{,}583 visits currently labelled TD:

\begin{table}[htbp]
\centering
\begin{tabular}{@{}lrrr@{}}
\toprule
\textbf{Arrived from} & \textbf{\arr\ TD} & \textbf{\arr\ Indet.} & \textbf{\arr\ PIGD} \\
\midrule
TD            & 0.849 & 0.065 & \textbf{0.086} \\
Indeterminate & 0.528 & 0.176 & 0.296 \\
PIGD          & 0.482 & 0.125 & \textbf{0.393} \\
\bottomrule
\end{tabular}
\caption{Where a patient came from changes the probability of moving to PIGD by
a factor of 4.6, even though they are in the same state now.}
\end{table}

\begin{keypoint}
\textbf{This is the signature of a latent state read through noise.} If a
patient's true phenotype is PIGD but one examination scores them TD, the next
examination is likely to score PIGD again --- producing exactly this dependence
on history. It is also consistent with the Weeks 3--4 finding that 35.2\% of
observed subtype changes reverse at the very next visit.

The consequence is that a chain fitted to the observed label is misspecified
\emph{whatever} its time parameterisation. That prediction is tested in
Section~5 and confirmed.
\end{keypoint}

% =============================================================================
\section{The three models}
% =============================================================================

\subsection{Discrete Markov chain --- the sketched baseline}

One transition matrix, applied per visit, interval ignored. Six free
parameters. This is the model the proposal describes as ``the probability of
remaining stable or transitioning at the next follow-up''.

\subsection{Continuous-time Markov chain}

A generator matrix $Q$ whose off-diagonal entries are instantaneous rates. The
probability of moving from state $i$ to $j$ over elapsed time $t$ is
\[
  P(t) = \exp(Qt),
\]
so every observed pair contributes at its own actual interval and the model can
be evaluated at any horizon. Fitted by maximising
$\sum_{\text{pairs}} \log P(\Delta t)[i, j]$ over six log-rates. Weeks 3--4
found transition probabilities scale strongly with the inter-visit gap, so this
was the expected improvement.

\subsection{Hidden Markov model --- latent state plus misclassification}

A latent phenotype $X(t)$ evolves as a continuous-time chain with generator
$Q$. What is recorded at a visit is not $X$ but a noisy reading $Y$, related by
a misclassification matrix
\[
  E[a, b] = P(\text{label recorded is } b \mid \text{true state is } a).
\]
The likelihood of a patient's whole sequence marginalises over every latent
path, computed by the forward algorithm. Fourteen free parameters: six rates,
six emission terms, two for the initial distribution.

This is the discrete-state counterpart of the proposal's stated observation
model, \emph{observed ratio = true ratio + measurement noise}.

\paragraph{An identifiability constraint.} Left unconstrained the
maximum-likelihood solution is degenerate, and identically so across all eight
random restarts: the latent Indeterminate state emits uniformly
($0.333/0.333/0.333$, i.e.\ its label carries no information) and PIGD becomes
absorbing at the parameter boundary. Requiring each state to be recorded
correctly more often than not ($E_{aa} \geq 0.5$) removes both. It costs
nothing: the held-out advantage over the discrete baseline is $+0.027$ with the
constraint against $+0.026$ without. A sweep over floors from $0.33$ to $0.80$
changes that advantage only between $+0.023$ and $+0.027$.

% =============================================================================
\section{Held-out results}
% =============================================================================

\begin{table}[htbp]
\centering
\begin{tabular}{@{}lrrr@{}}
\toprule
\textbf{Model} & \textbf{Log-lik / transition} & \textbf{Accuracy} & \textbf{Brier} \\
\midrule
HMM (latent + misclassification) & \textbf{$-0.6384$} & \textbf{0.769} & \textbf{0.3458} \\
Discrete Markov (per-visit)      & $-0.6653$ & 0.768 & 0.3625 \\
Continuous-time Markov           & $-0.6752$ & 0.766 & 0.3670 \\
Marginal (null)                  & $-0.8633$ & 0.625 & 0.5185 \\
\bottomrule
\end{tabular}
\caption{276 held-out patients, 1{,}002 transitions. Higher log-likelihood and
lower Brier are better.}
\end{table}

\subsection{The continuous-time chain does not help}

\begin{table}[htbp]
\centering
\begin{tabular}{@{}lrrrr@{}}
\toprule
\textbf{Interval} & \textbf{$n$} & \textbf{CTMC} & \textbf{DTMC} & \textbf{Difference} \\
\midrule
$\leq 1.5$ yr  & 890 & $-0.6534$ & $-0.6488$ & $-0.0046$ \\
$1.5$--$2.5$   &  71 & $-0.8151$ & $-0.7907$ & $-0.0244$ \\
$> 2.5$        &  41 & $-0.9054$ & $-0.8061$ & $-0.0993$ \\
\bottomrule
\end{tabular}
\caption{Held-out log-likelihood per transition by interval. The
continuous-time chain loses at every interval and by most at long ones ---
where respecting the gap should have helped most.}
\end{table}

\noindent
This is a negative result and worth stating plainly: the Weeks 3--4 finding
that transition probabilities scale with the interval was correct, but
correcting for it inside a chain fitted to the observed label does not improve
prediction. The reason is Section~3 --- the constraint $P(t) = \exp(Qt)$ forces
consistency across intervals that a noisily observed process does not have, so
imposing it on a misspecified model hurts rather than helps.

\subsection{Robustness}

\begin{table}[htbp]
\centering
\begin{tabular}{@{}lrr@{}}
\toprule
\textbf{Model} & \textbf{Mean} & \textbf{SD} \\
\midrule
HMM                    & $\mathbf{-0.6781}$ & 0.0251 \\
Discrete Markov        & $-0.7064$ & 0.0250 \\
Continuous-time Markov & $-0.7113$ & 0.0257 \\
Marginal (null)        & $-0.8764$ & 0.0187 \\
\bottomrule
\end{tabular}
\caption{Held-out log-likelihood per transition over 10 independent
patient-level splits, each refitting every model from scratch.}
\end{table}

\begin{keypoint}
The HMM beats the discrete chain in \textbf{10 of 10} splits, by
$+0.0284 \pm 0.0094$ log-likelihood per transition. Paired $t$-test
$p = 5.4 \times 10^{-6}$; Wilcoxon signed-rank $p = 0.002$. The ranking
HMM $>$ discrete $>$ continuous-time $>$ null does not change in any split.
\end{keypoint}

\begin{figure}[htbp]
\centering
\includegraphics[width=\textwidth]{model_comparison.png}
\caption{Every model on every split. Grey lines join the four models within one
split, so a line that never crosses shows a ranking that is consistent rather
than an artefact of averaging.}
\end{figure}

\clearpage

% =============================================================================
\section{What the fitted model says}
% =============================================================================

\subsection{Latent rates}

\begin{table}[htbp]
\centering
\begin{tabular}{@{}lrrr@{}}
\toprule
\textbf{From} & \textbf{\arr\ TD} & \textbf{\arr\ Indeterminate} & \textbf{\arr\ PIGD} \\
\midrule
TD            & $-0.0994$ & 0.0937 & \textbf{0.0057} \\
Indeterminate & 0.1156 & $-0.4713$ & 0.3557 \\
PIGD          & \textbf{0.0000} & 0.0605 & $-0.0605$ \\
\bottomrule
\end{tabular}
\caption{Latent generator $Q$, instantaneous rates per year.}
\end{table}

\begin{keypoint}
\textbf{Progression is sequential.} The direct TD~\arr~PIGD rate is 0.0057 per
year against 0.0937 for TD~\arr~Indeterminate --- \textbf{16 times smaller} ---
and the direct PIGD~\arr~TD rate is estimated at exactly zero. The model was
not told the states are ordered; it recovered that from the data.

This is what one should expect if the underlying quantity is the continuous
TD/PIGD ratio: to move from above 1.15 to below 0.90 a patient must pass
through the intermediate band. The same chain fitted to the observed label puts
the direct TD~\arr~PIGD rate at 0.0583, ten times higher, because apparent
direct jumps are misclassification.
\end{keypoint}

\subsection{How reliable is the clinical label?}

\begin{table}[htbp]
\centering
\begin{tabular}{@{}lrrr@{}}
\toprule
\textbf{True state} & \textbf{recorded TD} & \textbf{recorded Indet.} & \textbf{recorded PIGD} \\
\midrule
TD            & \textbf{0.941} & 0.038 & 0.021 \\
Indeterminate & 0.250 & \textbf{0.500} & 0.250 \\
PIGD          & 0.106 & 0.052 & \textbf{0.842} \\
\bottomrule
\end{tabular}
\caption{Estimated misclassification matrix $E$. The Indeterminate row sits at
the imposed floor of 0.50, so it should be read as ``at most half of
Indeterminate labels are correct'' rather than as a point estimate --- the
unconstrained likelihood drives it to $1/3$.}
\end{table}

\noindent
This corroborates a Weeks 3--4 result obtained by a completely different route:
changes \emph{into} Indeterminate were the ones that reversed most often
(56.0\% for TD~\arr~Indeterminate, 50.0\% for PIGD~\arr~Indeterminate). Two
independent analyses agree that the Indeterminate label is the least reliable
of the three.

\subsection{The phenotype is more stable than its label}

\begin{table}[htbp]
\centering
\begin{tabular}{@{}lrr@{}}
\toprule
\textbf{State} & \textbf{Observed label (CTMC)} & \textbf{Latent phenotype (HMM)} \\
\midrule
TD            & 3.41 & \textbf{10.06} \\
Indeterminate & 0.44 & \textbf{2.12} \\
PIGD          & 2.28 & \textbf{16.52} \\
\bottomrule
\end{tabular}
\caption{Mean years before leaving a state. Treating the label as the state
makes every phenotype look transient; separating measurement error from genuine
change recovers a much slower underlying process.}
\end{table}

\begin{figure}[htbp]
\centering
\includegraphics[width=0.72\textwidth]{sojourn_comparison.png}
\caption{The same comparison. Both quantities are model estimates, not raw
counts.}
\end{figure}

\clearpage

% =============================================================================
\section{The deliverable: probabilities of subtype transition}
% =============================================================================

\begin{table}[htbp]
\centering
\begin{tabular}{@{}llrrr@{}}
\toprule
\textbf{Current} & \textbf{Horizon} & \textbf{P(TD)} & \textbf{P(Indeterminate)} & \textbf{P(PIGD)} \\
\midrule
\textbf{TD} & now      & 1.000 & 0.000 & 0.000 \\
            & 1 year   & 0.910 & 0.071 & 0.019 \\
            & 2 years  & 0.834 & 0.111 & 0.055 \\
            & 3 years  & 0.769 & 0.133 & 0.099 \\
            & 5 years  & 0.660 & 0.149 & 0.191 \\
\midrule
\textbf{Indeterminate} & now     & 0.000 & 1.000 & 0.000 \\
            & 1 year   & 0.088 & 0.636 & 0.276 \\
            & 2 years  & 0.137 & 0.424 & 0.440 \\
            & 3 years  & 0.163 & 0.300 & 0.538 \\
            & 5 years  & 0.182 & 0.185 & \textbf{0.633} \\
\midrule
\textbf{PIGD} & now     & 0.000 & 0.000 & 1.000 \\
            & 1 year   & 0.003 & 0.047 & 0.950 \\
            & 2 years  & 0.009 & 0.075 & 0.916 \\
            & 3 years  & 0.018 & 0.091 & 0.891 \\
            & 5 years  & 0.036 & 0.107 & 0.857 \\
\bottomrule
\end{tabular}
\caption{Probability of occupying each subtype at each horizon, given the
current phenotype. This is the output format specified in the proposal.}
\end{table}

\begin{figure}[htbp]
\centering
\includegraphics[width=\textwidth]{progression_curves.png}
\caption{The same information continuously. Because the latent process runs in
continuous time, any horizon can be read off, not only the four tabulated.}
\end{figure}

\noindent
\textbf{The clinically useful line.} A patient currently Indeterminate has a
63\% chance of being PIGD within five years, against 19\% for a patient
currently TD. Indeterminate is not a mild or intermediate finding --- on these
estimates it is the strongest short-term predictor of conversion in the model.

\subsection{Calibration}

\begin{table}[H]
\centering
\begin{tabular}{@{}lrrr@{}}
\toprule
\textbf{Predicted-probability bin} & \textbf{$n$} & \textbf{Mean predicted} & \textbf{Observed rate} \\
\midrule
0.048--0.067 & 204 & 0.063 & 0.069 \\
0.067--0.083 & 207 & 0.075 & 0.101 \\
0.083--0.225 & 190 & 0.117 & 0.121 \\
0.225--0.643 & 202 & 0.445 & 0.391 \\
0.643--0.810 & 199 & 0.755 & 0.789 \\
\bottomrule
\end{tabular}
\caption{Held-out calibration for P(PIGD at next visit). Predictions track
observed frequencies across the range; the model is not systematically
over- or under-confident.}
\end{table}

% =============================================================================
\section{Limitations}
% =============================================================================

\begin{itemize}
  \item \textbf{The emission floor is an assumption, not an estimate.} The
        Indeterminate row rests on it. The honest statement is that at most
        half of Indeterminate labels are correct; the data alone would say
        one third.

  \item \textbf{Time-homogeneity is assumed and is violated.} Fitting the
        continuous-time chain separately to early and late follow-up gives a
        TD~\arr~PIGD rate of 0.045 per year before year 3 and 0.085 after ---
        nearly double. The model averages over disease duration.

  \item \textbf{No covariates.} The model uses only the subtype sequence. None
        of the baseline predictors assembled in Weeks 1--2 (RBDSQ, SCOPA-AUT,
        UPSIT, MoCA, DaTscan, age) enter it yet.

  \item \textbf{Three states, not a continuous ratio.} This is a discrete
        approximation to the proposal's design. 16.3\% of labelled visits have
        no finite ratio, which is one reason the discrete form was built first.

  \item \textbf{Label availability is not random.} Carried over from Weeks
        3--4: visits are lost preferentially for patients who cannot complete
        the pull test, who are more advanced. Progression estimates here remain
        lower bounds.

  \item \textbf{Held-out, not prospective.} Validation is a random split of one
        cohort. Nothing here has been tested on an external cohort.
\end{itemize}

% =============================================================================
\section{Anticipated questions}
% =============================================================================

\noindent\textbf{Q. The proposal asked for a Markov model. Why did you build three?}\\
Because the first two failed validation, and the failure was diagnosable. The
discrete chain is the model as sketched; the continuous-time chain was the
expected fix given the Weeks 3--4 interval finding; neither beat the other
meaningfully because both are misspecified in the same way. The Markov test in
Section~3 identified the reason, and the HMM is the model that reason points to.

\medskip
\noindent\textbf{Q. Isn't a 14-parameter model beating a 6-parameter one just overfitting?}\\
The comparison is on \emph{held-out patients} who were not used in fitting, and
it is repeated on 10 independent splits with a win in every one. Overfitting
would show as an advantage on training data that reverses out of sample. It
does not reverse.

\medskip
\noindent\textbf{Q. Why did the continuous-time model lose? Weeks 3--4 said the interval matters.}\\
The interval does matter --- that finding stands. But $P(t) = \exp(Qt)$ imposes
a strong consistency condition across intervals, and a process observed with
35\% immediate reversion does not satisfy it. Imposing a correct-looking
constraint on a misspecified model makes it worse, which is why the loss is
largest at long intervals.

\medskip
\noindent\textbf{Q. How do you know the latent states mean anything?}\\
Two independent checks. First, the model recovers an ordering it was never
given: direct TD~\arr~PIGD transitions are 16 times rarer than transitions via
Indeterminate, which is what a continuous underlying ratio implies. Second, the
estimated label reliability ranks Indeterminate lowest, matching the Weeks 3--4
reversion analysis computed by an entirely different method.

\medskip
\noindent\textbf{Q. Can this predict when a patient will convert?}\\
Partly. The table gives the probability of being PIGD at any horizon for a
patient in a given state, which answers ``how likely, by when''. It cannot yet
give a patient-specific answer, because no patient-level covariates or random
effects are in the model. That is the next stage.

\medskip
\noindent\textbf{Q. Is the model Markov in the latent state? Have you tested that?}\\
\textbf{No, that has not been tested.} The Markov property was tested and
rejected for the \emph{observed} label. The HMM assumes the \emph{latent}
process is Markov, and its better held-out performance is consistent with that
but does not establish it. A latent-state Markov test is a reasonable next
check.

% =============================================================================
\section{What comes next}
% =============================================================================

\begin{description}[leftmargin=1.5cm,style=nextline]
  \item[Patient-specific progression] The current model gives one set of rates
    for the whole cohort. The proposal's hierarchical structure --- a
    population trajectory plus subject-specific random effects --- is what
    turns this into an individualised prediction.
  \item[Move to the continuous ratio] The three-state model is a discrete
    approximation. The proposal's design treats the observed ratio as a noisy
    reading of a latent continuous phenotype, which subsumes this model and
    resolves the Indeterminate identifiability problem, since a threshold band
    is not a state.
  \item[Relax time-homogeneity] Rates are demonstrably not constant over
    disease duration; allowing them to depend on time since diagnosis is a
    contained extension.
  \item[Add baseline covariates] The malignant-phenotype markers are already
    assembled and 97--100\% complete at baseline. This is the second half of
    the proposal: does baseline diffuse-synucleinopathy burden predict
    progression, and progression \emph{stability}?
\end{description}

\vspace{1em}
\noindent\textbf{Reproducing this report.}
\begin{verbatim}
    python scripts/03_progression_model.py     # ~20 min: 10 splits + bootstrap
    python -m pytest tests/ -q                 # 118 tests
\end{verbatim}
\noindent
Stage 3 writes 14 CSV tables, 5 figures and \code{report.txt} to
\code{outputs/model/}. Every number quoted here appears there.

\end{document}
